% MA-SQLGrid Information reconstruction, round 3, 2026-09-13.
% Empirical evidence retained; formal properties are not new efficacy experiments.
\documentclass[information,article,submit,moreauthors]{Definitions/mdpi}
\firstpage{1}
\makeatletter\setcounter{page}{\@firstpage}\makeatother
\pubvolume{1}\issuenum{1}\articlenumber{0}\pubyear{2026}\copyrightyear{2026}
\datereceived{}\daterevised{}\dateaccepted{}\datepublished{}
\usepackage{booktabs}
\usepackage{algorithm}
\usepackage{algorithmic}
\setlength{\headheight}{20pt}
\setlength{\emergencystretch}{3em}
\newcommand{\masql}{MA-SQLGrid}
\AtBeginDocument{\pdfstringdefDisableCommands{\def\times{x}}}
% This is an unpublished revision: suppress the template-generated dummy DOI.
\makeatletter
\AtBeginDocument{\fancyfoot[R]{}\apptocmd{\ps@plain}{\fancyfoot[R]{}}{}{}}
\makeatother

\Title{MA-SQLGrid: An Auditable Coordination and Evaluation Framework for Power-Grid Text-to-SQL}
% ORCID: none declared for all authors on 24 August 2026; per the MDPI template, ORCID commands are omitted.
\Author{Bijing Liu $^{1,2}$, Chenglong Sun $^{1,2}$ and Yong Yang $^{1,2,}$*}
\AuthorNames{Bijing Liu, Chenglong Sun and Yong Yang}
\address{$^{1}$ \quad NARI Group Corporation (State Grid Electric Power Research Institute), Nanjing 211106, Jiangsu Province, China;\\
$^{2}$ \quad Beijing Kedong Electric Power Control System Co., Ltd., Beijing 100080, China}
\corres{Correspondence: yangyong1@sgepri.sgcc.com.cn}

\abstract{Executable SQL can still answer the wrong question. We investigate where reference-free validation loses correct candidates in a power-grid Text-to-SQL workflow. \masql{} records candidate production, eligibility, scoring, and tie resolution through five role interfaces and a deterministic controller. A finite-pool analysis separates unavailable answers, gate losses, scoring failures, and tie-order losses. In a development-visible synthetic GridDB case study with eight candidate slots for each of 180 questions, validation-aware and complete-witness selection matched 99 and 100 references, versus 76 for fixed-order selection and 129 for the strongest fixed source. The exhaustive post-result partition found 42 questions without a correct candidate, no correct-candidate gate losses, six scoring failures, and 33 or 32 tie-order losses. Both selectors had mixed-correctness top sets on 51 questions, despite slot ties on 130. Additional generation and BIRD Mini-Dev studies characterize context effects and non-grid portability; they do not replicate the core selection diagnosis across databases. The findings locate the limitations of the tested evidence: execution and invariance checks preserve available correct candidates but do not reliably distinguish them. They support a bounded diagnostic framework, not an end-to-end advantage from five-role coordination or independently validated power-grid semantics.}
\keyword{text-to-SQL; information systems; query validation; candidate selection; semantic ambiguity; multi-agent workflows; reproducibility}

\begin{document}

\section{Introduction}

A natural-language database interface must do more than produce SQL that runs. Consider a request for assets with scheduled maintenance after June. A query may filter a work-order status, a scheduled date, or a maintenance-log timestamp; each may be syntactically valid and return plausible rows. Successful execution settles none of these interpretations. This example is illustrative, not an additional benchmark item. It exposes the problem studied here: what evidence can a system use to choose between executable candidates without access to the reference answer?

Power-grid information systems make this distinction consequential because equipment, topology, measurements, and work orders encode different meanings for apparently similar fields. A numeric identifier may be categorical, and a timestamp may denote measurement, issuance, acknowledgement, or closure. The relevant database structure is therefore only one part of the question-to-query mapping. An inspectable interface should retain the interpretation, schema links, candidate statements, execution outcomes, and reason for selection. It should also make clear which of these records actually influenced the decision.

Spider, BIRD, and Spider~2.0 have advanced evaluation across schemas and enterprise workflows \cite{yu2018spider,li2023bird,lei2025spider2}. Schema-aware generation, constrained decoding, contextual retrieval, decomposition, and execution-guided repair address complementary obstacles \cite{wang2020ratsql,scholak2021picard,nan2023enhancing,pourreza2023dinsql,talaei2024chess}. Multi-role systems extend this design space through specialized generation and review stages \cite{lee2025team,wang2025macsql,shao2025qcmasql,xia2025r3}. However, a generator's ability to place a correct answer in a candidate pool and a selector's ability to recover it are different capabilities. The present study isolates this second problem alongside controlled generation and portability analyses.

We introduce \masql{}, a five-role interface coordinated through a shared blackboard. Intent analysis, schema mapping, candidate normalization, read-only validation, and named-state testing produce explicit records. The controller first checks eligibility, then ranks eligible candidates using question-derived evidence, and finally applies a disclosed tie rule. Reference answers enter only the offline evaluator after the decision has been sealed. This separation supports reproducible inspection of selection failures; it is not, by itself, evidence of better answers.

The central question is where a reference-free workflow loses the opportunity to return a correct candidate. Execution, output-shape, ordering, and state evidence can affect different stages, so a single final accuracy cannot identify the appropriate remedy. We examine three linked research questions:
\begin{enumerate}
\item Where do failures occur between candidate availability, eligibility, scoring, and tie resolution, and how do the selectors compare with every fixed source in the same historical pool?
\item Which recorded signals actually change selection, and which ambiguities remain under constructed-state checks and changes in candidate order?
\item How do context and generation choices shape the supporting GridDB and public non-grid results, and what limits do those studies place on the selection diagnosis?
\end{enumerate}

The contribution is a linked specification and diagnostic study. The framework first separates reference-free decision records from offline correctness evaluation. Elementary finite-pool properties then define an exhaustive failure partition, distinguishing absence of a correct answer from losses at eligibility, scoring, or tie resolution. Applying that partition to the frozen pool shows that no correct slot is gated out, while scoring and tie resolution leave available correct answers unselected. The complete fixed-source comparison and order audit expose the limits of the tested policy rather than claiming an accuracy advantage. The formal properties organize this diagnosis; they are not presented as a new general theory of SQL semantics.

Here, \emph{multi-agent} denotes five specialized role interfaces, some deterministic; the controller is not a sixth reasoning agent. In the historical-pool driver, schema grounding is recorded but does not affect selection, and the candidate adapter makes no model calls. GridDB is a development-visible synthetic case study, whereas BIRD tests non-grid portability. These boundaries rule out a causal claim that five-role coordination improves end-to-end accuracy. They leave a narrower, testable contribution: identifying where available evidence rejects candidates and where it fails to distinguish their meanings.

\section{Related Work}

\subsection{Schema Linking, Context Construction, and Evaluation}

Text-to-SQL systems map a natural-language request and database schema to an executable query. Spider established cross-domain evaluation on unseen databases, and BIRD expanded the setting toward larger databases with richer content \cite{yu2018spider,li2023bird}. Spider~2.0 and MultiSQL further emphasize enterprise-scale schemas, varied SQL operations, and context-dependent workflows \cite{lei2025spider2,li2024multisql}. These benchmarks make schema generalization visible, but their scores still depend on the database engine, evaluator, split, and model configuration.

Schema linking connects question spans to tables, columns, values, and foreign-key relations \cite{wang2020ratsql,lei2020reexamining,liu2022semantic}. Full-schema context protects recall but increases token cost and distraction. Compact context reduces the search space but can omit a bridge table or coded field. Retrieval, graph pathfinding, and bidirectional table--column selection seek a better balance \cite{bozdemir2026schema,liu2025solidsql,hao2025genlink,wang2025linkalign,safdarian2026schemagraphsql,nahid2026bidirectional}. Value examples add another signal because the intended field may be clearer from stored codes than from its name. The GridDB studies in this paper examine these context and value choices as observable components of a larger workflow.

Evaluation also needs more than one notion of match. String comparison can penalize equivalent SQL, while result equality can reward different queries that happen to coincide on one database snapshot. Distilled test suites evaluate denotation across multiple database states to expose some accidental equivalence \cite{zhong2020semantic}. \masql{} separately reports strict execution equality, constructed-state agreement, and bounded metamorphic invariance. These endpoints describe execution behavior under the stated transformations; expert-reviewed business semantics remain a separate validation layer.

The distinction between evaluation and selection is central here. Zhong et al.\ construct test suites with reference-query information to assess predicted SQL \cite{zhong2020semantic}. The historical-pool selector in this study cannot use reference correctness: its witnesses compare a candidate's behavior across specified transformations. The new analysis joins correctness only after the decision is sealed, to locate losses along the candidate--gate--score--tie sequence. It is therefore a diagnostic use of an evaluator, not a replacement for test-suite semantic evaluation or a claim to solve general SQL equivalence.

\subsection{Multi-Role Generation, Validation, and Selection}

Constrained decoding, decomposition, contextual grounding, and execution-guided repair address different stages of SQL production \cite{scholak2021picard,nan2023enhancing,sun2023sqlprompt,pourreza2023dinsql,talaei2024chess}. Recent workflows add question rewriting, attention-oriented decomposition, partition-and-drill reasoning, preference optimization, active database probing, or dialect-aware bootstrapping \cite{mao2024dartsql,xie2024deasql,luo2024ptdsql,zhai2025executionfeedback,xie2026sdesql,zhang2025exesql}. These methods mainly improve generation or repair. \masql{} focuses on the interfaces that connect candidate production to database-enforced validation and final selection.

Role specialization also appears in systems that train separate language models to store and retrieve SQL-based conversational memory \cite{lee2025team}; that task is distinct from generating answers to database queries. In Text-to-SQL, multi-agent systems have used decomposition and tools, question-dependent routing, and review--rebuttal--revision consensus \cite{wang2025macsql,shao2025qcmasql,xia2025r3}. \masql{} uses \emph{multi-agent} to denote five stable role interfaces recorded in a shared blackboard trace. The SQL Synthesizer may package externally produced candidates, whereas validation, state analysis, and final selection use explicit records and deterministic rules. Reproducibility therefore applies to these handoffs and the adjudication procedure, not to multi-model dialogue.

Read-only execution and evidence-aware adjudication are complementary. Database controls prevent an admissible evaluation from mutating the source, whereas adjudication compares only candidates that satisfy the required safety, execution, and state-evidence gates. A stored trace cannot guarantee sound reasoning, but it can reveal whether an error entered through intent analysis, schema mapping, generation, execution, or selection. This diagnostic role motivates the framework even when candidate generation is held fixed.

\subsection{Database Interpretation and Validation in Information Systems}

Two related studies in \emph{Information} illustrate why the object of evaluation must be explicit. Avignone et al.\ investigate natural-language descriptions of entity--relationship schemas, using construct-oriented hints and multiple description metrics \cite{avignone2025exploring}. Their task evaluates how a model explains database structure, rather than whether a selected SQL statement answers an operational question. Wardani et al.\ study AI-assisted SQL exercises and evaluate student answers against canonical solution variants \cite{wardani2026sql}. Their educational setting supplies task-specific answer definitions and user evidence that are absent from the present grid case study.

These studies motivate a distinction between structural interpretation, answer assessment, and application usefulness. In \masql{}, schema records support inspection, the frozen evaluator measures reference agreement, and neither substitutes for an independently reviewed business meaning. Consequently, the framework is compared with Text-to-SQL methods at the level of the problem addressed, while numerical comparisons are restricted to methods executed under the same disclosed protocol. No cross-paper accuracy ranking is inferred.

\subsection{Power-Grid Text-to-SQL}

Power-system data combine topology, equipment attributes, measurements, time series, costs, constraints, and maintenance records. RTS-GMLC and SimBench provide reproducible engineering structures that can be transformed into relational databases \cite{barrows2020rts,meinecke2020simbench}. Their availability does not by itself produce an expert-reviewed Text-to-SQL benchmark: intended projections, units, ordering, tie handling, and time semantics still require domain adjudication.

DKA-SQL is a direct power-grid Text-to-SQL precedent that adapts domain knowledge to grid questions \cite{bian2025dkasql}. Its corpus and evaluation boundary differ from the present study. \masql{} instead investigates a systems architecture for traceable role coordination, read-only validation, constructed-state evidence, and deterministic candidate selection. GridDB supplies a controlled grid case study, BIRD supplies public cross-database portability evidence, and RTS-GMLC/SimBench assets identify a path toward broader structural validation. Table~\ref{tab:literature-position} summarizes this positioning without treating results from different datasets and call budgets as directly comparable.

\begin{table}[htbp]
\caption{Positioning against adjacent Text-to-SQL approaches; entries describe reported design emphasis, not comparable performance.}
\label{tab:literature-position}\centering\small
\resizebox{\textwidth}{!}{%
\begin{tabular}{p{0.20\textwidth}p{0.22\textwidth}p{0.22\textwidth}p{0.26\textwidth}}
\toprule
Approach & Principal emphasis & Evaluation setting & Difference from \masql{}\\
\midrule
RGISQL \cite{li2024rgisql} & Grammatical and relational graph modeling & Public Text-to-SQL tasks & Learned generator rather than a traced executor/adjudicator boundary\\
Zero-shot prompting \cite{zhou2025zeroshot} & Prompt strategies and LLM behavior & Public benchmarks & Prompt comparison without the five-role blackboard\\
Schema retrieval \cite{bozdemir2026schema} & Embedding/vector-store context selection & LLM SQL generation & Retrieval quality rather than fail-closed named-state selection\\
DKA-SQL \cite{bian2025dkasql} & Domain knowledge adaptation & Power-grid Text-to-SQL & Complementary domain-adaptation focus; \masql{} studies traceable validation and candidate selection\\
\masql{} & Five-role trace, read-only validation, state evidence, and deterministic selection & GridDB case study plus BIRD portability & Integrates generation-to-selection evidence for power-grid Text-to-SQL\\
\bottomrule
\end{tabular}
}
\end{table}

\section{Materials and Methods}

\subsection{Problem Formulation and System Overview}

For question $x_i$, read-only database $D_i$, and introspected schema $S_i$, a generator produces one or more candidates
\begin{equation}
\mathcal{Y}_i=\{\widetilde{y}_{i1},\ldots,\widetilde{y}_{ik}\}.
\end{equation}
Each candidate must satisfy a single-statement read-only contract. The accepted leading form is \texttt{SELECT} or \texttt{WITH}; comments, multiple statements, write operations, schema changes, attachments, and pragmas are rejected. A candidate is eligible for adjudication only when it is safe and executable.

The offline evaluator may later compare the selected query with a reference result. Let $E_i=1$ when the selected SQL executes and its result equals the frozen reference result under the documented evaluator, and $E_i=0$ otherwise. Reference SQL, reference results, $E_i$, and any derivative correctness field are excluded from the inputs to all five roles and the deterministic controller. Evaluation occurs only after the blackboard is sealed.

This boundary yields three distinct states. A candidate is \emph{admissible} when its syntax and requested operations satisfy policy, \emph{executable} when it completes within the configured database limits, and \emph{correct under the frozen evaluator} when its result matches the reference after sealing. Admissibility is neither executability nor correctness. The manuscript preserves all three labels because collapsing them would let an executable but semantically wrong query appear validated.

The framework's response contract contains the selected SQL, stable candidate and source identifiers, referenced tables and columns, execution-state and result hashes, state-coverage summary, abstention or selection reason, and the sealed trace digest. Figure~\ref{fig:coordination} shows how these records move through the five-role architecture. The controller is not a sixth reasoning model: it schedules the roles, checks their typed messages, applies the deterministic rule, and seals the shared blackboard trace before offline evaluation.

\begin{figure}[htbp]
\centering
\includegraphics[width=\textwidth]{figures/fig_ma_sqlgrid_dual_panel.pdf}
\caption{\masql{} information flow and candidate lifecycle. \textbf{(a)} The Query Analyst and Schema Cartographer post question and schema records to the shared blackboard, and the SQL Synthesizer normalizes externally supplied candidates. The Validation Engine records read-only execution evidence, while the Metamorphic-State Critic records named-state evidence when required. \textbf{(b)} Hard eligibility gates precede the effective evidence score $e(y)=10I_{\mathrm{shape}}+5I_{\mathrm{order}}+5V$, $V\in[0,1]$, and the deterministic tie or abstention rule. Reference results are loaded only after the decision and trace digest are sealed.}
\label{fig:coordination}
\end{figure}

The response contract supports traceability, not user authorization. A calling application must authenticate the user, restrict the accessible database view, and decide whether returned rows may be displayed. Generated SQL and execution success do not replace those institutional controls.

\subsection{Read-Only Validation and Evidence-Aware Adjudication}

Eligibility always requires safety and successful execution. The archived 100-point representation assigns 40 points for each of those two properties, followed by 10 for question-derived aggregate-shape conformity, 5 for question-derived ordering conformity, and up to 5 for lexical value hits. Because every ranked candidate has already passed both hard gates, the first 80 points are constant and contain no ranking information. The effective evidence score is therefore $e(y)=10I_{\mathrm{shape}}(y)+5I_{\mathrm{order}}(y)+5V(y)$ for an eligible candidate $y$, with $V(y)\in[0,1]$. When named-state evidence is optional, the controller selects an $\arg\max e(y)$ and ties preserve the predefined candidate order. When named-state evidence is required, incomplete coverage or failure of the specified pass threshold makes $y$ ineligible before the same ranking. If the eligible set is empty, the controller abstains. This gate--score--tie function is mathematically equivalent to the archived implementation and shows that safety and execution are eligibility controls rather than discriminative evidence. The reported weight, pass-threshold, and order analyses are descriptive because the project history does not establish that every design choice preceded access to derived outcomes for the 180 items.

The 40/40/10/5/5 allocation, equivalently hard safety and execution gates followed by a 10/5/5 evidence score, is an illustrative deterministic policy rather than a learned or calibrated weighting. Its formal invariants are that unsafe or non-executable candidates are never selected; required named-state coverage is complete before threshold evaluation; gold-derived correctness is absent from the decision record; equal scores use the predefined order; and an empty eligible set yields abstention. These invariants support conformance testing and replay determinism, while semantic correctness and calibrated confidence require additional evidence.

Question-derived shape checks are deliberately weak. They detect cues such as ``how many'', ``highest'', ``first'', and explicit limits and compare them with the projected aggregate or ordering structure. They do not infer business semantics such as whether ``active'' refers to an asset, alarm, or work-order state. Lexical value hits likewise reward only presented values traceable to the question and context; they are not evidence that the associated column is the intended one. These low-bandwidth features keep adjudication inspectable but explain why large tie sets remain possible.

Abstention has two layers. Pre-ranking abstention occurs when no candidate passes safety and execution requirements or when required state coverage is incomplete. Post-ranking ambiguity is currently resolved by stable order because no margin threshold was frozen before outcome access in v3. The high tie rates observed later show that future work should freeze an ambiguity policy (for example, abstaining on unresolved top ties or requesting a user clarification) before examining new outcomes.

Robustness under tested conditions is represented by four observable mechanisms: database mutation denial, bounded execution, complete named-state eligibility, and deterministic failure handling. Candidate-order stability is assessed separately in Results. Algorithm~\ref{alg:coordination} states the resulting reference-isolated coordination rule.

\begin{algorithm}[htbp]
\caption{Reference-isolated deterministic coordination}
\label{alg:coordination}
\begin{algorithmic}[1]
\REQUIRE question $x_i$, schema $S_i$, external candidates $\mathcal{Y}_i$, read-only executor, named state evidence
\STATE post structured intent and schema grounding to the append-only blackboard
\STATE canonicalize and deduplicate externally produced candidates
\FOR{each candidate $y\in\mathcal{Y}_i$}
  \STATE apply single-statement read-only validation
  \STATE record execution and available reference-free state evidence
\ENDFOR
\IF{required named-state coverage is incomplete for a candidate}
  \STATE mark that candidate ineligible (fail closed)
\ENDIF
\IF{no safe executable candidate meets all specified gates}
  \STATE abstain
\ELSE
  \STATE apply the fixed validation, state-evidence, coverage, and order rule
\ENDIF
\STATE post the decision and seal the blackboard
\STATE load reference results only for offline evaluation
\end{algorithmic}
\end{algorithm}

\subsection{Formal Scope of Evidence-Based Selection}
\label{sec:formal-scope}

The following properties state what follows from the disclosed rule and what does not. They assume a finite candidate set, correctly implemented eligibility checks, fixed evidence records, and a deterministic score. They do not assume that the available features recover the user's intended semantics.

For question $i$, let $\mathcal{A}_i\subseteq\mathcal{Y}_i$ be the candidates passing all required gates, and define the top set
\begin{equation}
\mathcal{T}_i=\{y\in\mathcal{A}_i:e(y)=\max_{z\in\mathcal{A}_i}e(z)\}.
\label{eq:topset}
\end{equation}
For nonempty $\mathcal{A}_i$, the current controller returns the first member of $\mathcal{T}_i$ under the fixed order $\pi$; otherwise it returns $\bot$ (abstention). Let $c_i(y)\in\{0,1\}$ denote correctness under the frozen offline evaluator, with $c_i(\bot)=0$. This function is used only to analyze a sealed decision.

\begin{Proposition}[Eligibility and finite-pool bounds]
\label{prop:pool-bound}
The controller selects no candidate outside $\mathcal{A}_i$. For any selector $f_i$ constrained to $\mathcal{A}_i\cup\{\bot\}$,
\begin{equation}
c_i(f_i)\leq \max_{y\in\mathcal{A}_i}c_i(y)
\leq \max_{y\in\mathcal{Y}_i}c_i(y),
\label{eq:poolbound}
\end{equation}
where the maximum over an empty set is defined as zero.
\end{Proposition}
\begin{proof}
The gate is applied before ranking, so a returned candidate belongs to $\mathcal{A}_i$. Its binary correctness cannot exceed the largest correctness value in that set. Inclusion of $\mathcal{A}_i$ in $\mathcal{Y}_i$ gives the second inequality; abstention contributes zero.
\end{proof}

Averaging Equation~\ref{eq:poolbound} over questions separates three sources of error: a correct candidate may be absent, excluded by a gate, or available but not selected. The maxima require reference outcomes and are oracle diagnostics, not deployable policies. A best fixed source is a different comparator: it uses one source for every question and need not attain the question-wise oracle.

For diagnosis, a failed decision can be assigned to exactly one of four stages, examined in order: no correct candidate in $\mathcal{Y}_i$; correct candidates exist but none survives in $\mathcal{A}_i$; an eligible correct candidate exists but none reaches $\mathcal{T}_i$; or the top set contains a correct candidate that the tie rule does not select. Together with a correct selected answer, these categories partition the evaluated questions. The partition locates a failure in the disclosed decision rule; it does not identify the business-semantic cause of an incorrect candidate or imply that a different tie policy can attain every question-wise oracle simultaneously.

\begin{Proposition}[Order dependence of the top set]
\label{prop:ties}
If $\mathcal{T}_i$ contains one candidate, changing candidate order does not change the selection. If it contains multiple distinct candidates, each can be selected by an order placing it before the other top candidates. Correctness varies with order if and only if the top set contains both evaluator-correct and evaluator-incorrect candidates.
\end{Proposition}
\begin{proof}
Only candidates in $\mathcal{T}_i$ maximize the fixed score. A singleton has no tie to resolve. For a non-singleton, placing a given member first makes it the selected maximizer. The possible correctness values are therefore exactly $\{c_i(y):y\in\mathcal{T}_i\}$.
\end{proof}

This is a per-question property over available orders. A single global ordering constrains all questions jointly, so independently attainable per-question maxima cannot simply be summed to predict a globally attainable score. Likewise, different SQL strings can return the same result. Slot ties, distinct-string ties, and mixed-correctness ties must not be treated as interchangeable counts.

\begin{Proposition}[Limits of observational evidence]
\label{prop:indistinguishability}
Suppose two admissible interpretations of a request yield identical inputs to a deterministic selector, including the question, candidates, and recorded evidence, but have disjoint sets of correct candidates. A selector returning the same candidate on those inputs cannot be correct under both interpretations. Abstention avoids choosing between them but does not resolve the request.
\end{Proposition}
\begin{proof}
Identical inputs produce the same deterministic output. No candidate belongs to both correct sets by assumption. The returned candidate is therefore incorrect under at least one interpretation; $\bot$ provides no candidate answer.
\end{proof}

The assumption is relevant when a business definition is missing rather than encoded in the available evidence. For example, status-based and date-based readings of \emph{scheduled} can remain distinct even when both queries execute, return the same width, and pass storage-preserving witnesses. This is an illustrative ambiguity, not a claim that every observed tie has this cause. Additional schema semantics, a discriminating test, or user clarification may resolve it; changing a source-order tie breaker alone supplies no missing meaning.

These properties make a concrete prediction about the scope of validation, not its success rate: a check that passes every surviving candidate contributes no discrimination within that set. Constant eligibility points cannot affect ranking. A new witness is useful for selection only when it changes eligibility or provides a ranking feature; whether that change is beneficial remains an empirical question. The historical-pool and component results below test these distinctions without treating the proofs as evidence of improved accuracy.

\subsection{Five-Role Architecture and Coordination}

\masql{} comprises five specialist role interfaces coordinated by the controller shown in Figure~\ref{fig:coordination}. Messages receive consecutive sequence numbers, and the complete trace is serialized canonically to a SHA-256 digest. Each role has restricted inputs and produces one typed output for the next stage.

The \textbf{Query Analyst} produces a structured question-only intent record containing lexical tokens, detected aggregations, ordering requirements, and an explicit limit when present. It does not select SQL or see correctness. The \textbf{Schema Cartographer} maps lexical intent to tables, columns, and foreign-key edges. The current skeleton uses deterministic overlap and records unmatched tokens; a richer GridDB adapter adds disclosed domain rules and graph expansion.

The \textbf{SQL Synthesizer} packages external SQL without model calls, canonicalizes formatting, removes exact duplicates while preserving first occurrence, and assigns stable identifiers. The \textbf{Validation Engine} used by the historical-pool study opens the database with \texttt{mode=ro\&immutable=1}, enables \texttt{query\_only}, disables extensions, and installs an authorizer. It denies writes, DDL, transactions, attachments, pragmas, metadata enumeration, dangerous functions, and reads outside table/column allowlists. Progress and fetch handlers enforce time, opcode, and row limits. Every attempt retains database/SQL hashes, failures, timing, row count, and result hash without hidden retry.

A later executor revision adds function allowlisting, projected-width, raw-cell-byte, and total-result-byte limits. Fourteen tests cover BLOB denial, total-budget and boundary behavior, trace retention, oversized text or results, and wide projections. These controls postdate the frozen historical candidate and evidence ledgers; they therefore strengthen the implementation boundary without changing those stored artifacts. Table and column allowlists constrain query scope but do not authenticate users or provide row-level authorization; both executor variants establish tested SQLite mutation denial and scoped execution bounds.

The threat model assumes untrusted SQL but a trusted host, SQLite library, configured allowlist, and immutable database. It does not address a compromised process, OS escape, side channels, built-in malicious extensions, or upstream authorization errors. Production use requires process isolation, least-privilege views, authenticated context, logging, and independent security review.

Failure handling is intentionally fail-closed and does not retry. A parse denial, authorizer denial, timeout, opcode limit, row limit, width limit, cell limit, or total-result limit creates an attempt record and makes that candidate ineligible for the current decision. The controller cannot silently weaken a limit or substitute a repaired query. A repair procedure may be studied as a separately specified generation step, as in BIRD B3, but its extra physical call and visibility must remain explicit.

The \textbf{Metamorphic-State Critic} (historically named the Counterfactual Critic) consumes named state results, rejects duplicate state identifiers, and records evaluated, passed, failed, and missing states. When state evidence is optional, missing evidence remains diagnostic and does not alter the previously defined no-state behavior. When the caller requires state evidence, coverage must equal the complete specified state set and the pass threshold must be met; otherwise the candidate is ineligible. An incomplete 1/1 record therefore cannot outrank a complete 10/11 record. Gold-relative labels from the constructed-state study remain forbidden during selection.

The blackboard is append-only at the application layer: a role posts a typed message and cannot rewrite an earlier message. The controller checks role, message type, sequence number, and required fields before acceptance. Canonical serialization fixes key order and whitespace before hashing. This makes two runs comparable at the trace level and exposes an unexpected hidden field or changed handoff. It does not prevent a malicious host process from altering memory or files; process isolation and remote attestation are outside scope.

Information flow is intentionally asymmetric. The Analyst receives the question but no candidate SQL. The Cartographer receives the question, intent record, and introspected schema. The Synthesizer receives external candidate strings and assigns identities but cannot execute them. The Validator receives one candidate plus a scoped executor. The Critic receives named-state outcomes but no reference correctness. Only the deterministic controller receives the eligible evidence summaries. Gold enters a separate offline evaluator after the decision message and board digest exist.

\subsection{Data Sources, Visibility, and Evaluation Roles}

GridDB-Maintenance-v2 v0.1 is a synthetic SQLite case study with eight tables and 98 rows: \texttt{asset\_types} (6), \texttt{assets} (18), \texttt{grid\_topology} (9), \texttt{locations} (8), \texttt{maintenance\_logs} (8), \texttt{sensor\_readings} (26), \texttt{technicians} (8), and \texttt{work\_orders} (15). It contains 200 natural-language--SQL records: 20 development records and 180 factorial evaluation records. The evaluation records include 66 easy, 91 medium, and 23 hard items and map to 70 normalized gold-SQL structural clusters, of which 58 are singletons. The evaluation partition had been visible during project development; it is a controlled finite-set case study, not a sealed production benchmark.

BIRD Mini-Dev contributes 500 public items over 11 SQLite databases. The protocol fixed item order, prompts, two local models, evaluator, Python 3.10.11/SQLite 3.40.1, and zero retries. Each backbone made 2500 calls and produced 2000 final predictions; independent re-execution found zero score mismatches. Two incomplete predecessor runs were excluded from every scientific denominator and retained only in the supplementary execution record.

The GridDB reliability release contains 18 states: the canonical snapshot, three insertion permutations, and 14 additional schema-valid stress states. The scorer produced 25,920 prediction-state rows ($1440\times18$). A prediction-blinded automated protocol admitted a 66-question subset for the primary 15-state analysis and retained 114 order-sensitive questions as diagnostics. These constructed states are not operator-certified grid snapshots.

The evaluation distinguishes inherited predictions, deterministic recomputations, new protocol-bound runs, and diagnostic implementation records. These labels preserve chronology without treating software replay as an accuracy experiment. Table~\ref{tab:resources} summarizes the three resources that enter the main empirical argument. RTS-GMLC and SimBench remain potential external-validation resources because their current question--SQL pairs have not completed qualified-human semantic review.

\begin{table}[htbp]
\caption{Resources used in the main empirical argument. Each resource answers a different research question and retains its own denominator.}
\label{tab:resources}
\centering
\small
\setlength{\tabcolsep}{2pt}
\begin{tabular}{p{2.6cm}p{4.5cm}p{3.0cm}p{3.0cm}}
\toprule
Resource & Role in the study & Sample or unit & Primary endpoint\\
\midrule
GridDB v0.1 & Controlled grid-schema generation, component, and historical-pool studies & 1 database; 8 tables; 98 rows; 20 development + 180 evaluation questions & Strict execution equality, component contrasts, and reference match after selection\\
Constructed-state suite & Stability of retained predictions under specified database transformations & 18 states; 1440 predictions; 66-question primary subset & Fifteen-state logical-AND agreement and witness coverage\\
BIRD Mini-Dev & Portability of the prompting and execution workflow to public non-grid databases & 500 items across 11 SQLite databases & Execution accuracy and database-grouped method contrasts\\
\bottomrule
\end{tabular}
\end{table}

\subsection{Experimental Protocols and Evaluation Boundaries}

Table~\ref{tab:protocol-master} maps each protocol to its unit, endpoint, and evidence boundary. The studies answer complementary questions and are interpreted separately because their candidate sources, call budgets, and visibility histories differ.

\begin{table}[htbp]
\caption{Experimental design and inferential role of each protocol. Composition intervals describe sensitivity to the observed finite-corpus grouping, not population confidence.}
\label{tab:protocol-master}\centering\scriptsize
\setlength{\tabcolsep}{3pt}
\begin{tabular}{p{2.3cm}p{2.6cm}p{3.2cm}p{4.3cm}}
\toprule
Study & Unit/groups & Primary endpoint & Research role\\
\midrule
GridDB factorial & 180 questions/cell; 70 normalized-SQL groups & Execution equality with nine-test Holm correction & Context and hint effects for two model snapshots\\
Component E1/E2 & 170 paired or 180 questions; 61 or 70 groups & Execution equality and grouped composition intervals & Presented-value and multi-candidate-selection effects\\
Constructed states & 66 questions; 12 structural groups & Fifteen-state logical AND; exact $2^{12}$ signs & Agreement under prediction-blind transformations\\
BIRD & 500 items; 11 databases & Execution accuracy and 12 Holm-adjusted contrasts & Public non-grid portability of the workflow\\
Coordination replay & 180 questions & Pool coverage and complete finite-pool counts & Trace and handoff reproducibility\\
Historical-pool selection & 180 questions; shared 5760-attempt ledger & Reference match, rescue/harm, ties, and order sensitivity & Evidence-aware selection within the observed pool\\
\bottomrule
\end{tabular}
\end{table}

GridDB factorial outcomes are inherited and development-visible; the component procedure was fixed before its calls on development-visible items; constructed states were prepared without reading predictions; and BIRD used a fixed public protocol. The historical-pool study sealed every decision before loading raw reference outcomes, but same-item derived outcomes had been accessed earlier in development. Its selector comparisons are consequently descriptive, although their execution is deterministic and traceable.

\subsection{Controlled Context and Hint Study}

The paired design crosses context package $c\in\{0,1\}$ and composite hint $h\in\{0,1\}$. F00 uses full DDL and global values without the hint; F01 adds the hint; F10 uses compact domain-grounded context without the hint; and F11 combines compact context and the hint. The same 180 questions appear in all four cells for each backbone, yielding 720 Qwen and 720 Granite predictions.

Because each question appears in every prompt cell, cell differences are paired within question. The two backbone snapshots form parallel finite-corpus experiments, not independent replications drawn from a model population. Difficulty labels and normalized-SQL clusters describe the corpus structure. Cluster-based randomization reduces the influence of near-duplicate SQL forms, while composition-sensitivity intervals show how estimates vary when observed clusters are resampled. Neither procedure converts this development-visible synthetic set into a population sample.

Qwen2.5-Coder-7B Q4\_K\_M and Granite-3.3-8B Q4\_K\_M were run at temperature zero with fixed snapshots. Each prompt produced exactly one response. The parser retained the first SQL candidate through the first semicolon; the validator then applied the read-only policy. There was no candidate ranking, multi-agent negotiation, or repair in this experiment.

Strict execution equality is primary. A common-target projected-column indicator is a secondary manipulation check. Paired finite-set contrasts use normalized-SQL clusters as dependence proxies, documented sign-flip procedures, composition-sensitivity bootstrap intervals, and Holm correction \cite{cameron2008bootstrap,canay2017randomization,holm1979simple}. The intervals describe sensitivity to the observed corpus composition, not population confidence.

\subsection{Component, Constructed-State, and Public-Database Studies}

The separately specified component study contains 700 scored calls. E1 compares candidate generation with and without presented value evidence on 170 paired questions in 61 groups. E2 requests up to three candidates and compares a deterministic reference-free selector with the first candidate on 180 questions in 70 groups. Question-level paired differences are averaged with equal question weights; bootstrap resampling and sign assignments operate at group level. Selection is sealed before gold loading. Rescue, harm, and oracle-at-three are descriptive; oracle-at-three is a gold-only upper-bound diagnostic. The three two-member Holm families are E1 across backbones, E2 across backbones, and the two cross-backbone modifiers.

The constructed-state reliability protocol reproduced all 1440 T0 labels and then executed every prediction over 18 states. The primary 66-question analysis uses a question-weighted logical-AND endpoint over 15 semantic states. A prediction passes only when it agrees with the reference on every included state. The nine specified contrasts form a separate Holm family; with only 12 structural groups, all $2^{12}=4096$ group-sign assignments are enumerated. The sharp-null interpretation assumes exchangeability of those 12 group signs. The small group count limits randomization resolution, and the groups are dependence proxies rather than sampled deployment clusters.

The BIRD protocol compares direct prompting, decomposition, schema selection, and bounded execution repair. Each method has a fixed call pattern and zero retries. Accuracy and method differences retain equal item weighting over 500 items; pairwise randomization assigns one common sign to all paired item differences within each of 11 databases and applies Holm correction across 12 comparisons. Its sharp-null interpretation assumes database-level sign exchangeability. With only eleven clusters, the attainable evidence is coarse, and the databases do not represent a random sample of power-grid systems.

The four protocols answer different questions and are not pooled. GridDB factorial cells test prompt-package manipulations on a controlled grid schema. E1 and E2 test presented values and reference-free selection. The constructed-state study tests agreement across transformed databases, and BIRD tests non-grid workflow portability.

\subsection{Fixed-Candidate Design for Isolating Selection}

The replay verifies SHA-256 for two 720-row prediction ledgers and the 25,920-row constructed-state ledger before reading them. For each question it pools the eight existing outputs from two backbones and four prompt cells, canonicalizes SQL, and removes duplicates. T0 fields provide execution and metadata-shape evidence. The Critic records named-state evidence as unavailable because existing state-agreement labels are gold-relative.

This preliminary replay is a coverage diagnostic. The candidates were generated under different models and prompt packages, and no selected output from this step is used as a performance estimate.

The historical-pool selection study reuses exactly eight source slots for each of the 180 GridDB evaluation questions: four Qwen prompt cells followed by four Granite prompt cells. It performs no LLM call and creates no new question or SQL label. Before selection, all three methods share one precomputed eight-by-four execution ledger: T0 plus three metamorphic witnesses for every slot. The shared ledger contains 5760 physical evidence-collection attempts; this is not a separate runtime cost assigned to each selector. The fixed-order control ignores the ranking evidence and chooses slot 0. Validation-only selection ranks candidates without named-state evidence, whereas complete-witness selection requires all three reference-free metamorphic outcomes. Incomplete coverage or no passing candidate causes abstention.

The three query-blind witnesses have distinct SHA-256 values and are built from T0 without reading questions, prompts, predictions, scores, or gold. M1 adds an irrelevant relation and rows; M2 adds harmless indexes and rebuilds storage; M3 adds nullable probe columns to three business tables. M1 and M2 preserve original-schema denotation. M3 preserves explicitly projected original columns but deliberately allows wildcard projection to fail. Query order is preserved when candidate SQL contains \texttt{ORDER BY}; otherwise rows are compared as a multiset. The selector sees only question text extracted from retained prompt ledgers, the schema, eight SQL strings, executor traces, and T0-to-witness equivalence. It does not see gold SQL, answer-shape metadata, correctness, difficulty, feature tags, required literals, or gold-relative constructed-state labels.

These witnesses test metamorphic relations rather than alternative ground-truth answers. M1 asks whether an irrelevant table changes behavior, M2 whether logically neutral storage changes alter the result, and M3 whether projection is stable when nullable columns are appended. A pass means that the observed result satisfies the specified relation; it does not mean the original query expresses the user's intent. A fail can reveal brittleness, as with wildcard projection, but may also reflect a transformation outside the intended query equivalence class. For this reason, witness outcomes enter eligibility only under the disclosed construction and remain reference-free diagnostics.

Candidate order is a first-class experimental variable because the controller uses stable order for equal scores. Original order places four Qwen slots before four Granite slots and follows the prompt-cell sequence within each backbone. A post-review order audit changes only this tie breaker after the same evidence has been collected and exactly enumerates all $8!=40{,}320$ global slot orders. It is a direct diagnostic of unresolved ranking information, not an alternative model or a new sample. The observed range motivates an abstention or clarification policy for future work.

Each selector used the same 1440 source slots and 5760 attempts, and no failed candidate was replaced or retried. The complete software-test matrix, protocol chronology, manifest hashes, and excluded predecessor records are provided in Supplementary Sections S1--S3.

Descriptive endpoints are strict execution accuracy over all 180 questions, coverage and abstention, complete-three-witness invariance of the selected candidate, and rescue/harm relative to C000. The unified-evaluator order audit reports the exact correct-count distribution across all global slot orders. A descriptive risk--coverage curve covers questions whose top-score tie set is no larger than $k$, for $k=1,\ldots,8$, and uses the frozen original order within covered ties. Because same-item outcomes were accessible before these diagnostics were finalized, neither the permutation audit nor the risk--coverage curve is a confirmatory analysis.

A post-review reconciliation freezes evaluator \nolinkurl{MA-SQLGrid-GridDB-T0-shape-denotation-v1} and applies it to all eight fixed slots and the sealed selector choices. Candidate and gold SQL execute under the same bounded read-only T0 snapshot. Correctness requires successful execution, the frozen expected column count, and duplicate-preserving row equality (ordered when specified; floating absolute tolerance $10^{-6}$). Empty--empty rows are compared only after both shape gates pass. SQL canonicalization is used for artifact identity, not correctness. This reconciliation is post-result evidence rather than preregistration.

\subsection{Statistical Analysis}

The statistical unit is the question, and every question receives equal weight. The principal estimand is the mean paired difference on the observed finite corpus. Normalized-SQL groups for GridDB and databases for BIRD are dependence proxies used for grouped resampling and sign assignment; they are not random samples from a deployment population. Composition-sensitivity intervals describe how a finite-corpus estimate changes when the observed groups are reweighted and are not population confidence intervals.

The randomization procedures assume exchangeability of group-level signs under the tested sharp null. Exact enumeration is used for the 12-group constructed-state analysis ($2^{12}=4096$ assignments), while the larger documented families use the stated Monte Carlo procedures and seeds. All tests are two-sided, zero paired differences remain in the estimand and contribute no sign, and the Holm families are defined separately for the nine GridDB factorial effects, projected-column manipulation checks, three two-member component families, nine constructed-state contrasts, and 12 BIRD pairings. Count differences, paired effects, rescue/harm, and sensitivity results remain visible whether or not a Holm-adjusted decision is met.

The historical-pool comparison is reported through complete counts, paired rescue/harm, coverage, abstention, tie multiplicity, and candidate-order sensitivity. Because same-item outcomes were accessible during earlier development, these quantities describe the observed pool and do not estimate a population or causal multi-agent effect.

Appendix~\ref{sec:historical-question-statistics} retains the historical question-level reconciliation: paired question-composition bootstrap intervals and exact discordant-sign tests against C000, with Holm adjustment over nine unique comparisons. Unlike the grouped analyses above, these tests do not protect against dependence between questions sharing a SQL structure. Their small values must not be read as confirmation across templates or databases. The main selection table therefore reports finite-pool counts, differences, and paired rescue/harm without inferential labels. The new stage decomposition joins frozen candidate correctness to recorded eligibility, scores, and sealed choices only after selection; it produces no online feature and is explicitly post-result.

\subsection{Implementation and Reproducibility}

The local generation experiments used the identified Qwen2.5-Coder-7B-Instruct Q4\_K\_M and Granite-3.3-8B-Instruct Q4\_K\_M snapshots under frozen protocols. The BIRD formal protocol fixed Python 3.10.11, SQLite 3.40.1, model order, prompts, evaluator, and zero retries. The five-role controller, read-only executor, state-evidence gate, and deterministic replay have regression and adversarial tests. Main-text results always retain the control version actually used to generate them; later executor corrections are reported as implementation evidence and are not applied retroactively.

Detailed computational records, artifact hashes, test inventories, execution history, and rights information are provided in the Supplementary Materials.

\subsubsection{Generative-AI Assistance and Author Responsibility}

The local SQL-generation experiments used the identified Qwen2.5-Coder-7B-Instruct Q4\_K\_M and Granite-3.3-8B-Instruct Q4\_K\_M snapshots under the protocols described above. During manuscript preparation, the authors also used the OpenAI Codex agent interface for prose editing, experimental suggestions, Python/\LaTeX{} code review, computational verification, and local build and visual-quality checks. The exact backend model identifier and version were not retained for every session.

All executable suggestions were run locally against retained inputs, and every reported value was obtained from retained or independently recomputed records rather than generated prose. The authors verified the analyses, interpretations, citations, declarations, and final text and retain full responsibility for the manuscript. Codex output was not treated as qualified power-system expert annotation or independent adjudication. External machine-generated question--SQL candidates remain labeled machine silver. All six scientific figures are deterministic code-native graphics; no generative-image output is used as quantitative evidence.

\section{Results}

\subsection{Evidence Summary and Reading Guide}

The experiments produced three main observations. First, schema context, structural hints, and presented values affected the two model snapshots differently; no single context strategy dominated across every analysis. Second, under the unified evaluator, historical-pool reference matches were 76/180 for C000, 99/180 for validation ranking, and 100/180 with complete witness evidence, while the fixed Qwen F01 source reached 129/180. Third, both evidence-aware selectors produced top-score ties on 130/180 questions and remained sensitive to candidate order. Validation evidence changed choices relative to C000, but the available signals often remained too coarse for stable ranking and did not outperform the best fixed source.

The historical-pool comparison and exhaustive failure partition address RQ1. Role-utilization, tie, and constructed-state diagnostics address RQ2. Generation, presented-value, and BIRD studies address RQ3 as supporting evidence about candidate production and evaluation boundaries. BIRD does not rerun the historical-pool selector, so its database diversity cannot be credited as external replication of RQ1. Implementation conformance is summarized after the empirical results and documented fully in the Supplementary Materials.

\subsection{Validation-Aware Selection in the Historical Candidate Pool}

The replay covered all 180 GridDB questions without new generation. At least two distinct SQL candidates were available for 173 questions, and 172 had at least two eligible candidates after safety and execution checks. Seven questions had only one unique candidate, and one had fewer than two eligible candidates. These counts describe interface coverage; the full candidate-multiplicity distribution is reported in the numerical supplement.

\subsubsection{Reference Matches, Coverage, and Rescue/Harm}

The deterministic descriptive re-execution retained 5760 candidate--state attempts; 332 were non-executable or authorization/resource/SQL failures and remained in the all-attempt ledger. Every question nevertheless had at least one eligible candidate, so both adjudicating selectors covered 180/180 questions and abstained on 0/180. This zero-abstention result is a consequence of at least one eligible slot plus forced stable-order resolution of post-ranking ties, not confidence-calibrated coverage. Table~\ref{tab:offline} reports the recorded rules.

Keeping all 332 failures matters for two reasons. Removing them would overstate the executability of the historical pool, and retaining their slot identities verifies that selectors did not receive replacements or retries. Full question-level coverage means only that another slot was eligible for every question; it does not mean every candidate was valid or that an operational system should always answer.

\begin{table}[htbp]
\caption{Descriptive unified-evaluator comparison for the historical eight-slot pool. Differences and paired rescue/harm use C000 as baseline. C000 and Qwen F00 contain identical normalized SQL. All 180 questions are retained; no population or cross-template inference is attached to these counts. Historical question-level intervals and tests are preserved in Appendix~\ref{sec:historical-question-statistics}.}
\label{tab:offline}
\centering
\scriptsize
\resizebox{\textwidth}{!}{%
\begin{tabular}{lrrrr}
\toprule
Method & Correct/180 & Accuracy & Difference & Rescue/harm\\
\midrule
C000 / Qwen F00 & 76/180 & 0.4222 & 0 & 0/0\\
Qwen F01 & 129/180 & 0.7167 & +0.2944 & 56/3\\
Qwen F10 & 78/180 & 0.4333 & +0.0111 & 3/1\\
Qwen F11 & 108/180 & 0.6000 & +0.1778 & 46/14\\
Granite F00 & 77/180 & 0.4278 & +0.0056 & 6/5\\
Granite F01 & 100/180 & 0.5556 & +0.1333 & 45/21\\
Granite F10 & 74/180 & 0.4111 & -0.0111 & 6/8\\
Granite F11 & 108/180 & 0.6000 & +0.1778 & 42/10\\
Validation-only selector & 99/180 & 0.5500 & +0.1278 & 25/2\\
Complete-witness selector & 100/180 & 0.5556 & +0.1333 & 26/2\\
\bottomrule
\end{tabular}}
\end{table}

Qwen F00 and historical C000 contain the same 180 SQL artifacts. Their previously reported 76-versus-80 discrepancy came from evaluator policy: the historical row-only comparator treated Q104, Q107, Q110, and Q140 as matches because both results were empty, whereas all four candidate projections had the wrong column count. The unified evaluator applies the frozen shape gate before empty-row equality and reproduces all 1440 canonical fixed-slot verdicts. Under that evaluator, validation-only and complete-witness selection differ by one match and remain 30 and 29 matches below Qwen F01, respectively.

Figure~\ref{fig:offline-diagnostics} juxtaposes reference matches and witness invariance. The one-match selector increment focuses interpretation on validation ranking and tie policy, not a broad witness benefit.

\begin{figure}[htbp]
\centering
\includegraphics[width=0.90\textwidth]{figures/fig05_offline_selector_diagnostics.pdf}
\caption{Selector endpoints over 180 historical-pool questions. Reference matches are loaded only after board sealing; invariance denotes agreement across the three constructed witnesses.}
\label{fig:offline-diagnostics}
\end{figure}

Both evidence-aware selectors produced multiple highest-scoring candidates on 130/180 questions, with mean top-tie sizes of 5.4000 and 5.3889. With original-order tie breaking, 178/180 selections came from Qwen slots, reflecting source order rather than learned preference.

Most apparent source preference is mechanically induced. When all eight candidates tie, the first slot is selected regardless of backbone quality; when seven tie, the earliest surviving slot has the same advantage. The observed 178 Qwen selections measure how often the available reference-free features failed to separate candidates before the stable-order rule acted, not a learned preference for one backbone.

The numerical supplement reports the complete multiplicity and selected-slot distributions, including all eight-candidate ties.

Validation-only ranking added 23 net reference matches relative to C000 through 25 rescues and 2 harms. Complete-witness selection changed the chosen candidate on only Q039, yielding one further match (100/180 versus 99/180), one additional rescue, no additional harm, and no coverage or abstention change. Table~\ref{tab:q039} gives the constructed trace for this sole selector difference. Neither selector exceeded the fixed Qwen F01 source at 129/180.

\subsection{Tie Structure, Order Sensitivity, and the Q039 Case}

\begin{table}[htbp]
\caption{Q039 is the sole validation-only versus complete-witness selection difference. ``Match'' denotes agreement with the synthetic reference.}
\label{tab:q039}
\centering
\footnotesize
\setlength{\tabcolsep}{3pt}
\begin{tabular}{p{2.2cm}p{4.0cm}p{4.0cm}p{1.8cm}}
\toprule
Field & Validation-only & Complete witnesses & Interpretation\\
\midrule
Request & \multicolumn{2}{p{8.8cm}}{``Show the first scheduled work order after June 2024.''} & Status/date ambiguity\\
Projection & \texttt{SELECT *} & \texttt{SELECT work\_order\_id, scheduled\_date} & Wildcard vs. explicit fields\\
Shared filter & \multicolumn{2}{p{8.8cm}}{\texttt{scheduled\_date > '2024-06-01' ORDER BY scheduled\_date ASC LIMIT 1}} & Ambiguous time scope\\
M3 outcome & Wildcard width changes & Two-column projection preserved & Projection stability\\
Reference result & No match & Match & One-question trace\\
\bottomrule
\end{tabular}
\end{table}

M3 added nullable columns to \texttt{work\_orders}; consequently, complete-witness eligibility rejected the wildcard projection and retained the explicit \texttt{work\_order\_id, scheduled\_date} projection. The full SQL strings are supplied in the numerical supplement. Neither query resolves the request's status and date-boundary ambiguity through qualified review, so the extra reference match is interpreted as a projection-stability trace rather than an engineering correctness judgment.

Table~\ref{tab:sensitivity} replaces the non-comparable historical row-only sensitivity values with an exact audit under the unified evaluator. The audit enumerated all $8!=40{,}320$ global slot orders while preserving each question's frozen evidence scores. Validation-only selection ranged from 95 to 128 matches and complete-witness selection also ranged from 95 to 128; both had a median of 107. The original orders produced 99 and 100 matches, respectively, while the reverse order produced 116 for both selectors. Even the outcome-exposed maximum remained below the fixed Qwen F01 source at 129/180.

These permutations are diagnostics, not candidate policies: no order is preferred after outcomes are observed. A second descriptive audit abstained whenever the top-score set was not unique. It covered only 50/180 questions and matched 24 of them (0.4800), so top-tie size alone was not a calibrated confidence signal. The range across global orders and the poor strict-abstention result motivate a semantic tie or abstention rule frozen before evaluation on an untouched pool.

A transparent post-run hash audit emitted no SQL text and found duplicate normalized SQL in 154/180 pools. Mean validation and complete-witness top sets fell from 5.400 and 5.389 slots to 3.150 and 3.139 unique strings; the mean pool contained 5.322 unique strings. Deduplication reduces apparent multiplicity but does not remove ambiguity. The outcome-aware right-step tie-size AURC values, 0.4826 and 0.4795, are diagnostics rather than calibration results.

\begin{table}[htbp]
\caption{Exact global candidate-order sensitivity under the unified shape-and-denotation evaluator. All 40,320 orders are outcome-aware diagnostics; original and reverse are two members of that enumeration.}
\label{tab:sensitivity}
\centering
\small
\resizebox{\textwidth}{!}{%
\begin{tabular}{lrrrrrr}
\toprule
Selector & Original & Reverse & Minimum & Median & Mean & Maximum\\
\midrule
Validation-only & 99 & 116 & 95 & 107 & 109.845 & 128\\
Complete witness & 100 & 116 & 95 & 107 & 109.929 & 128\\
\bottomrule
\end{tabular}}
\end{table}

The post-review utilization audit then traced which role outputs actually enter the historical-pool selector. Query Analyst fields feed aggregation, order, and lexical-match validation features. By contrast, the Schema Cartographer's grounding is posted to the blackboard but is not passed to candidate generation, validation, the critic, or adjudication in this frozen driver. The candidate provider is an adapter for stored SQL and makes no model calls. Table~\ref{tab:role-ablation} therefore reports score-component diagnostics rather than claiming a prospective removal of five independently active agents. Safety and executability gates remain active in every variant.

\begin{table}[htbp]
\caption{Post-review role-utilization and single-component diagnostics under the unified evaluator. Changes, gains, and losses are relative to the stated parent and use the original candidate order.}
\label{tab:role-ablation}
\centering
\small
\resizebox{\textwidth}{!}{%
\begin{tabular}{llrrrr}
\toprule
Selector variant & Parent & Correct/180 & Choices changed & Gains & Losses\\
\midrule
Validation original & --- & 99 & 0 & 0 & 0\\
Without all Query-Analyst-derived score terms & Validation original & 76 & 64 & 2 & 25\\
Without shape term & Validation original & 99 & 0 & 0 & 0\\
Without order term & Validation original & 77 & 49 & 1 & 23\\
Without lexical value-hit term & Validation original & 99 & 14 & 1 & 1\\
Complete witness original & --- & 100 & 0 & 0 & 0\\
Without constructed-state gate & Complete witness & 99 & 1 & 0 & 1\\
Without recorded-only schema grounding & Complete witness & 100 & 0 & 0 & 0\\
\bottomrule
\end{tabular}}
\end{table}

The order feature accounts for most of the observed post-review change in this decomposition: removing it changes 49 validation choices and lowers the match count by 22. Removing lexical value hits changes 14 choices but exchanges one gain for one loss, while removing the shape term changes none. Removing constructed-state eligibility changes only Q039. These same-item, outcome-aware diagnostics localize behavior in the frozen pool; they do not estimate the benefit of an independently implemented role on unseen questions.

\begin{table}[htbp]
\caption{Exhaustive post-result failure partition under the unified evaluator. Each question appears in exactly one row for each selector. Correctness means reference agreement, not expert-reviewed business semantics.}
\label{tab:failure-partition}\centering\small
\begin{tabular}{lrr}
\toprule
Decision outcome & Validation-only & Complete witness\\
\midrule
No correct candidate in the original pool & 42 & 42\\
All correct candidates removed by eligibility gates & 0 & 0\\
Eligible correct candidates below the highest score & 6 & 6\\
Correct top candidate lost through tie order & 33 & 32\\
Selected candidate correct & 99 & 100\\
\midrule
Total questions & 180 & 180\\
\bottomrule
\end{tabular}
\end{table}

Table~\ref{tab:failure-partition} connects the formal bounds to the recorded decisions. A correct candidate exists in 138 pools and survives eligibility in every one; neither selector removes any evaluator-correct source slot in this dataset. Six questions lose their remaining opportunity at scoring. Both selectors have mixed-correctness top sets on 51 questions, of which the original order resolves 33 and 32 incorrectly. Thus 130 slot ties do not mean 130 correctness-sensitive ambiguities. These counts distinguish candidate availability, gate retention, score discrimination, and final tie resolution without treating any oracle choice as a usable policy.

An automated evaluator-state taxonomy further separates execution failure, output-shape mismatch, and wrong denotation after both shape gates pass. C000 has 1 execution error, 85 shape mismatches, and 18 wrong denotations among its 104 failures. Validation-only selection has 0, 42, and 39, respectively; complete-witness selection has 0, 41, and 39. The best fixed Qwen F01 source has 5 execution errors, 1 shape mismatch, and 45 wrong denotations. Thus the frozen selector removes many inadmissible or shape-inconsistent candidates but leaves a larger share of its remaining failures beyond the evaluator's shape gate. This automatic taxonomy does not distinguish wrong status meaning, units, time boundaries, topology intent, or professional usefulness; those categories remain subject to independent expert review.

The numerical supplement contains all eight GridDB cells, 12 component endpoints, eight 15-state cells, the exact unified-evaluator order histogram, the descriptive risk--coverage curve, the 360-row item-level tie ledger, selected-origin counts, and the Q039 trace. The main-text findings use different questions, candidates, call patterns, and endpoints, so they are not combined into a single framework accuracy. Across these studies, the common diagnostic finding is that context and validation evidence can change observable behavior, while stable semantic selection remains unresolved.

\subsection{Context and Component Effects on SQL Generation}

All 1440 scheduled predictions are present, and independent SQLite re-execution reproduced the stored execution verdicts. Table~\ref{tab:cells} reports strict execution equality, and Figure~\ref{fig:cells} visualizes the same audited cell estimates.

\begin{table}[htbp]
\caption{Strict execution equality over 180 GridDB questions per cell.}
\label{tab:cells}
\centering
\small
\setlength{\tabcolsep}{4pt}
\begin{tabular}{lcccc}
\toprule
Backbone & F00 full/no hint & F01 full/hint & F10 compact/no hint & F11 compact/hint\\
\midrule
Qwen & 76/180 (0.4222) & 129/180 (0.7167) & 78/180 (0.4333) & 108/180 (0.6000)\\
Granite & 77/180 (0.4278) & 100/180 (0.5556) & 74/180 (0.4111) & 108/180 (0.6000)\\
\bottomrule
\end{tabular}
\end{table}

\begin{figure}[htbp]
\centering
\includegraphics[width=0.94\textwidth]{figures/results/fig01_v2_cells.pdf}
\caption{Audited GridDB cell point estimates. Each cell contains all 180 attempts; the figure is descriptive and does not replace the documented factorial contrasts.}
\label{fig:cells}
\end{figure}

The common-target projected-column diagnostic is 0.5222 in F00 for both backbones. Qwen obtains 0.9667, 0.5889, and 0.9611 in F01, F10, and F11; Granite obtains 0.8778, 0.5667, and 0.9222. These values show that returning the requested number of columns is easier than returning the intended rows. In Qwen F01, projected-column conformity is 0.9667 while execution equality is 0.7167.

The eight cells also show that no single prompt summarizes the design. Adding the hint to full context increases the point estimate for both backbones, whereas adding compact context without the hint leaves execution equality close to F00. Under compact context, the hint produces 108 correct answers for both models, but this numerical convergence arises from different unhinted baselines and does not establish identical behavior. Reporting every cell preserves these interactions instead of selecting only the highest value.

The 0.2500 gap between Qwen F01 projected-column conformity and strict execution equality warns against structural proxy endpoints. A model may recognize that one column, an aggregate, or an ordered row is requested while still choosing the wrong table, join, predicate, value, or boundary. Shape diagnostics can localize errors and inform adjudication, but they cannot replace result equality or qualified semantic review.

\subsubsection{Factorial Effects and Multiplicity}

Table~\ref{tab:factorial-nine} reports the complete primary family. For Qwen, the composite-hint execution effect is +0.2306, the package--hint interaction is $-0.1278$, and the compact-package main effect is $-0.0528$. Granite's corresponding estimates are +0.1583, +0.0611, and +0.0139. None of the nine primary execution tests survives Holm correction. Qwen's hint effect has raw/adjusted $p=0.01372/0.09604$, its interaction $0.00919/0.07352$, and the cross-backbone hint modifier $0.00640/0.05760$. Accordingly, no primary factorial execution effect or modifier is described as statistically nonzero.

\begin{table}[htbp]
\caption{Complete primary factorial family for strict execution equality. Estimates are question-weighted finite-set differences; intervals are 95\% normalized-SQL-group composition-sensitivity intervals. Each sign-randomization used 70 groups, 100,000 draws, and the displayed raw values were Holm-adjusted jointly over all nine rows.}
\label{tab:factorial-nine}\centering\scriptsize
\resizebox{\textwidth}{!}{%
\begin{tabular}{llrrrr}
\toprule
Scope & Effect & Estimate & Composition interval & Raw $p$ & Holm $p$\\
\midrule
Qwen & Compact-package main & -0.0528 & [-0.1131, 0.0000] & 0.10523 & 0.52614\\
Qwen & Structural-hint main & 0.2306 & [0.0548, 0.4355] & 0.01372 & 0.09604\\
Qwen & Package $\times$ hint & -0.1278 & [-0.2454, -0.0341] & 0.00919 & 0.07352\\
Granite & Compact-package main & 0.0139 & [-0.0584, 0.0798] & 0.85021 & 1.00000\\
Granite & Structural-hint main & 0.1583 & [-0.0193, 0.3662] & 0.15480 & 0.61919\\
Granite & Package $\times$ hint & 0.0611 & [-0.0539, 0.2013] & 0.55890 & 1.00000\\
Granite--Qwen & Backbone $\times$ package & 0.0667 & [-0.0360, 0.1814] & 0.41478 & 1.00000\\
Granite--Qwen & Backbone $\times$ hint & -0.0722 & [-0.1229, -0.0276] & 0.00640 & 0.05760\\
Granite--Qwen & Three-way interaction & 0.1889 & [0.0090, 0.4321] & 0.07010 & 0.42060\\
\bottomrule
\end{tabular}}
\end{table}

The normalized-SQL mapping contains 70 groups, 58 singletons, and a maximum group size of 19. Question-level weighting is retained; groups are resampling and sign-assignment proxies, not verified authoring templates or sampled clusters. The randomization interpretation requires exchangeability of group signs under the tested sharp null, while the composition interval asks how the finite-corpus mean changes when those observed groups are reweighted. Neither assumption follows from the data, and the high singleton fraction limits protection against unmeasured dependence.

The displayed composition intervals are pointwise sensitivity intervals; they are neither simultaneous/familywise intervals nor inversions of the nine Holm-adjusted randomization tests. An interval that excludes zero can therefore coexist with a Holm non-rejection. The multiplicity decision, not interval exclusion, controls statements about the documented nine-test family.

No minimum practically important execution-accuracy difference, equivalence margin, or non-inferiority margin was specified before analysis. The study therefore cannot infer practical equivalence from a non-rejection, and the 1/180 complete-witness difference is a mechanism trace rather than evidence of a meaningful benefit.

In the separately corrected secondary family, the common-target hint effects survive: +0.4083 for Qwen (adjusted $p=0.000090$) and +0.3556 for Granite (adjusted $p=0.01944$). Because the hint supplies structural information later scored by this endpoint, these are manipulation-check results, not independent semantic evidence.

The multiplicity result changes the emphasis of the factorial study. Large raw contrasts motivate hypotheses about structural instructions, but none of the nine primary execution comparisons meets the documented familywise rule. The projected-column findings verify that the manipulation changed a closely aligned surface property. They do not show that the compact package or hint produces a backbone-independent accuracy gain. Future work should separate schema pruning, value presentation, normalization, and structural hints rather than retaining composite packages.

\subsubsection{Presented Values and Candidate Selection Components}

The component run completed all 700 scored calls under a prospectively frozen call-and-selection procedure on development-visible GridDB items. On the paired 170-question subset, Qwen first-candidate correctness increased from 83/170 to 101/170, a +0.1059 difference (95\% composition-sensitivity interval [+0.0282, +0.2013], adjusted $p=0.0310$). Granite remained 69/170 in both paired conditions, with an interval spanning zero. The cross-backbone modifier did not meet the frozen decision rule, so the Qwen result is not replicated across backbones and is not a five-role result.

The deterministic selector changed 24 Qwen and 36 Granite choices. Relative to the first candidate, it rescued eight and harmed one Qwen question, and rescued ten while harming none for Granite. Execution effects were +0.0389 and +0.0556; neither survived its two-test Holm family. Oracle-at-three values of 0.6389 and 0.4667 are gold-only diagnostics and are not deployable selector results.

For Qwen, the value-evidence intervention adds 18 correct first candidates when the 170 eligible V0 questions are compared with their V1 counterparts under the documented frozen component protocol; the corresponding Granite estimate is zero. The disagreement is more informative than averaging the backbones: the effect depends on how a frozen model uses schema and value context. Oracle-at-three measures candidate-pool headroom, whereas the reference-free selector measures how much of that headroom the recorded rule can recover without gold.

Complete component counts appear in the numerical supplement. On all 180 V1 questions, Qwen first/selector/oracle-at-three counts were 105/112/115, and Granite counts were 71/81/84; the paired 170-question intervention counts and all documented effects appear in Table~\ref{tab:component-effects}. These are component, not five-role, outcomes.

\begin{table}[htbp]
\caption{Complete component-effect families. Intervals are pointwise group-composition sensitivity summaries, not simultaneous intervals and not inversions of the Holm tests. Each family contains the two displayed members indicated by its identifier.}
\label{tab:component-effects}\centering\scriptsize
\resizebox{\textwidth}{!}{%
\begin{tabular}{llrrrrrr}
\toprule
Family & Contrast & $N$ & Groups & Estimate & Composition interval & Raw $p$ & Holm $p$\\
\midrule
E1 & Qwen V1--V0 first candidate & 170 & 61 & 0.1059 & [0.0282, 0.2013] & 0.01552 & 0.03104\\
E1 & Granite V1--V0 first candidate & 170 & 61 & 0.0000 & [-0.1902, 0.1705] & 1.00000 & 1.00000\\
E2 & Qwen selector--first V1 & 180 & 70 & 0.0389 & [-0.0081, 0.1071] & 0.50080 & 0.50080\\
E2 & Granite selector--first V1 & 180 & 70 & 0.0556 & [0.0075, 0.1232] & 0.06425 & 0.12850\\
Cross-E & Granite--Qwen E1 effects & 170 & 61 & -0.1059 & [-0.2701, 0.0394] & 0.26986 & 0.53971\\
Cross-E & Granite--Qwen E2 effects & 180 & 70 & 0.0167 & [-0.0062, 0.0457] & 0.37477 & 0.53971\\
\bottomrule
\end{tabular}}
\end{table}

The three two-member Holm families are E1 across backbones, E2 across backbones, and Cross-E across modifiers. The Granite E2 pointwise interval excludes zero while its adjusted decision does not meet the frozen decision rule; these summaries are not contradictory. The 61- and 70-group mappings retain question weights but are dependence proxies, not a probability sample. Figure~\ref{fig:components} visualizes these effects without replacing the adjusted decisions.

\begin{figure}[htbp]
\centering
\includegraphics[width=0.92\textwidth]{figures/results/figure_01_primary_effects.pdf}
\caption{Component effects under a prospectively frozen call-and-selection procedure on development-visible GridDB items. Intervals describe composition sensitivity; the frozen decision rule also requires multiplicity-adjusted evidence.}
\label{fig:components}
\end{figure}

\subsection{Constructed-State Stability}

The v5 scorer produced all 25,920 protocol-specified rows, and its T0 slice reproduced 1440/1440 snapshot labels. Table~\ref{tab:multistatecounts} and Figure~\ref{fig:multistate} report the 66-question automatic subset. The 15-state logical-AND rates range from 0.6212 (Granite F01) to 0.8182 (Granite F00). All nine Holm-adjusted values equal 1.0000 under exact cluster-sign enumeration. No context, hint, interaction, or cross-backbone multi-state effect meets the declared rule. These rates characterize constructed-state agreement, not population accuracy or human-certified equivalence.

The ordering under this endpoint differs from the single-snapshot table. Granite F00 has the highest logical-AND rate on the 66-question automatic subset even though Granite F11 has the larger T0 execution count over all 180 questions. The denominators and endpoints differ, so this is not a contradiction. The reversal shows why a robustness statement must name both the admitted question subset and the required state set.

\begin{table}[htbp]
\caption{Complete 15-state logical-AND counts for the prediction-blinded 66-question automatic subset.}
\label{tab:multistatecounts}
\centering
\small
\begin{tabular}{lcccc}
\toprule
Backbone & F00 & F01 & F10 & F11\\
\midrule
Qwen & 49/66 (0.7424) & 48/66 (0.7273) & 53/66 (0.8030) & 43/66 (0.6515)\\
Granite & 54/66 (0.8182) & 41/66 (0.6212) & 51/66 (0.7727) & 49/66 (0.7424)\\
\bottomrule
\end{tabular}
\end{table}

\begin{figure}[htbp]
\centering
\includegraphics[width=0.92\textwidth]{figures/results/fig04_semantic_reliability.pdf}
\caption{GridDB multi-state logical-AND agreement for the prediction-blinded automatic subset. The endpoint is a constructed-state robustness diagnostic.}
\label{fig:multistate}
\end{figure}

\subsection{Cross-Database Portability on BIRD}

Table~\ref{tab:bird} reports all eight cells. BIRD Mini-Dev is a public non-grid portability benchmark and does not validate power-grid semantics. B0--B2 use one physical generation call per item; B3 uses two, so the inherited comparison estimates frozen workflow efficacy under unequal call counts rather than a call-matched repair effect. Each backbone completed 2500 calls, 2000 final predictions, zero retries, and zero final-ledger omissions. Latency and token totals are unavailable in the retained summary and were not reconstructed.

\begin{table}[htbp]
\caption{Inherited BIRD Mini-Dev v1.1 results (non-grid portability evidence; 11 database clusters). Intervals are database-cluster composition-sensitivity intervals.}
\label{tab:bird}
\centering
\small
\setlength{\tabcolsep}{3pt}
\begin{tabular}{llrrrrp{1.5cm}}
\toprule
Backbone & Method & Correct/500 & Accuracy & Interval & Calls/item & Final-ledger omissions\\
\midrule
Qwen & B0 direct & 189/500 & 0.378 & [0.286, 0.470] & 1 & 0\\
Qwen & B1 decomposition & 151/500 & 0.302 & [0.203, 0.401] & 1 & 0\\
Qwen & B2 schema selection & 197/500 & 0.394 & [0.279, 0.511] & 1 & 0\\
\cmidrule(lr){1-7}
Qwen & B3 execution repair & 174/500 & 0.348 & [0.250, 0.452] & 2 & 0\\
Granite & B0 direct & 102/500 & 0.204 & [0.125, 0.284] & 1 & 0\\
Granite & B1 decomposition & 105/500 & 0.210 & [0.131, 0.294] & 1 & 0\\
Granite & B2 schema selection & 101/500 & 0.202 & [0.137, 0.273] & 1 & 0\\
\cmidrule(lr){1-7}
Granite & B3 execution repair & 118/500 & 0.236 & [0.159, 0.318] & 2 & 0\\
\bottomrule
\end{tabular}
\end{table}

After Holm adjustment, Qwen decomposition is 0.076 below direct prompting (adjusted $p=0.0430$), and schema selection is 0.092 above decomposition (adjusted $p=0.0117$). No Granite contrast survives correction. Different method orders across the two model snapshots argue against a universal prompting recipe. Two incomplete predecessor runs were retained for traceability and excluded from all reported denominators; their disposition is documented in the Supplementary Materials.

The BIRD results also constrain the meaning of ``agentic'' improvement. Decomposition reduces Qwen accuracy relative to direct prompting, while bounded repair uses twice the physical calls without becoming the best Qwen method. Granite shows a different descriptive ordering and no adjusted contrast. More stages or calls do not guarantee a better outcome. A future comparison must count physical generations, retries, context tokens, and failed attempts, not only final predictions.

\subsection{Implementation Reliability}

The five-role controller and executor were tested for typed handoffs, mutation denial, bounded execution, trace completeness, deterministic replay, and fail-closed state coverage. Failed attempts remain in the evidence ledger and cannot be silently replaced. The executor used to create the historical candidate and evidence ledgers enforced SQLite read-only mode, \texttt{query\_only}, authorization callbacks, disabled extensions, and time, opcode, and row bounds. A later revision added cell-byte, total-result-byte, width, and function controls; these tests strengthen the implementation boundary but do not reinterpret the frozen ledgers. The unified evaluator subsequently scored the stored fixed and selected SQL under one explicit shape-and-denotation rule. The historical selector ledger contains 3960 blackboard messages and 5760 database attempts (5428 executable and 332 SQLite failures) across 953 distinct SQL hashes; it contains no selector-stage model calls. Recorded local attempt time totals 839 ms, but timer granularity places both the median and 95th percentile at 0 ms, so these values are an audit of the ledger rather than a deployment-latency estimate. Supplementary Tables S1--S3 provide the complete test inventory and protocol attribution.

\section{Discussion}

\subsection{Validation and Selection Solve Different Problems}

The historical-pool result distinguishes a useful filter from a successful selector. Validation-aware ranking recovers 25 questions and harms two relative to C000, producing 23 additional reference matches. Yet the fixed Qwen F01 source achieves 30 more matches than validation-aware selection. The strongest-source comparison changes the interpretation: the pool already contains a better fixed policy, but the recorded evidence does not recover its performance. Proposition~\ref{prop:pool-bound} explains why candidate availability and selection ability must be reported separately.

The failure partition makes the redesign question more precise. Changing the selector alone cannot repair the 42 questions with no correct candidate. Improving eligibility retention is not the observed bottleneck because no correct slot was removed. The six scoring failures require evidence or a score that promotes a surviving correct candidate, whereas the 33 validation-only tie losses concern alternatives already at the maximum score. This is a diagnostic allocation of the observed errors, not a promise that the four interventions can be optimized independently. A future study can test whether the same allocation recurs across candidate pools before choosing which interface to redesign.

The automatic error taxonomy supports a more specific account than a general claim of improved reasoning. Relative to C000, the selectors reduce execution and shape failures while leaving many wrong denotations after those gates pass. They are therefore effective at some forms of rejection within this pool, but the remaining ranking problem concerns distinctions the current features represent only weakly. This is not proof that execution evidence is unhelpful in other systems. It shows why a validation metric should be interpreted through the errors it can detect, rather than equated with answer correctness.

The same distinction applies to the five-role architecture. Typed records identify where evidence enters the workflow and allow a decision to be replayed. They do not establish that every role contributes to prediction quality. In particular, recorded-only schema grounding cannot explain a change in historical-pool selection. The role-utilization audit is consequently as important as the architecture diagram: it separates an implemented interface from an active decision mechanism.

\subsection{Why the Tested Witnesses Add Little Discrimination}

Complete-witness eligibility changes only Q039, raising the count from 99 to 100. The witnesses were designed to test irrelevant-table, storage/index, and nullable-column transformations. These interventions probe specific invariances; they do not introduce the missing business definitions needed to interpret coded statuses or time fields. Limited incremental selection is therefore compatible with the narrow information supplied by the tests.

Q039 provides a concrete example. Nullable schema extension distinguishes a wildcard projection from an explicit two-column projection, while the intended meaning of scheduled maintenance remains unresolved. Agreement with the synthetic reference is the measured outcome; the trace is not independent expert adjudication of that reference. This example also prevents an overgeneralization of Proposition~\ref{prop:indistinguishability}: additional evidence can resolve a particular distinction, but a witness that resolves projection width need not resolve business intent.

For future systems, witness design should begin with a stated ambiguity and an expected discriminating relation. Merely increasing the number of state checks can add cost without separating candidates. Conversely, a gate may exclude a correct query if the transformation violates the query's intended equivalence class. Witness coverage, rejection patterns, and retained correct candidates should therefore be examined together. The present study does not establish an optimal witness set.

\subsection{Candidate Order Is Part of the Decision Rule}

Both selectors have top-score ties on 130 of 180 questions. Across all global slot orders, their reference-match counts range from 95 to 128, while reversal raises the original counts to 116. Proposition~\ref{prop:ties} identifies the mechanism: when the score cannot distinguish top candidates, the order supplies the final choice. The order is thus a substantive policy, not an innocuous implementation detail.

The empirical quantities need to remain distinct. Slot multiplicity includes duplicate SQL; deduplication reduces the average top set to about 3.1 unique strings but does not eliminate ambiguity. A tie between equivalent queries need not affect correctness, whereas a mixed-correctness tie can. The global-order enumeration captures their joint consequences for this finite pool; it does not identify a policy to select after observing outcomes.

Nor does uniqueness establish confidence. The strict no-tie rule covers 50 questions and matches only 24. A unique maximizer can still be wrong because the features themselves are imperfect. The relevant next experiment is therefore not simply to force agreement or remove ties. It is to compare prospectively frozen semantic evidence and clarification policies using both answer quality and coverage on untouched questions.

\subsection{Context Effects Depend on the Model and Evaluation Task}

The generation studies provide context for the selection failure. Structural hints strongly alter projected-column conformity, yet none of the nine primary GridDB execution contrasts survives its Holm correction. Presented values improve the Qwen first-candidate count on the paired component subset but leave Granite unchanged. BIRD likewise yields different method orderings: schema selection has the highest Qwen point estimate, while bounded repair has the highest Granite estimate and uses an additional call.

These observations are consistent with a trade-off between supplying useful context and introducing irrelevant or misleading context \cite{bozdemir2026schema,liu2025solidsql,hao2025genlink}. They do not show that one prompting strategy generalizes across backbones. The schema-description study of Avignone et al.\ similarly separates recognition of constructs from the quality of the resulting text \cite{avignone2025exploring}; here, the analogous distinction is between adopting the requested output shape and returning the intended rows.

A modular interface makes these alternatives easier to inspect, but modularity is not the treatment tested in the generation experiments. A causal architecture comparison would require matched models, physical calls, tokens, candidate counts, and evaluator conditions. The present studies instead establish the narrower fact that context and evidence choices can have different effects at different stages.

\subsection{Implications for Auditable Information Systems}

Three design implications follow from the specification and observed failures. First, separate policy admissibility, execution success, reference agreement, and independently reviewed business meaning in both logs and reports. Second, preserve candidate provenance, active feature values, gate failures, and tie rules so that a selected answer can be traced to its actual decision mechanism. Third, test whether added evidence distinguishes the surviving alternatives rather than rewarding checks that all of them pass.

These are recommendations motivated by this study, not measured deployment benefits. The framework does not authenticate a user, enforce institutional row-level policy, or establish that displayed records are appropriate for consequential decisions. Those responsibilities remain with the calling application and qualified reviewers. Similarly, trace hashing supports comparison of retained records under a trusted host; it does not make an untrusted host tamper-proof.

Figure~\ref{fig:evidence-map} links the empirical streams to these boundaries. Generation experiments test candidate production, state checks test specified invariances, and the historical-pool study tests selection. Keeping those endpoints separate makes the negative selection result informative to other information-system researchers without implying that this small grid case study measures a general failure rate.

\begin{figure}[htbp]
\centering
\includegraphics[width=0.88\textwidth]{figures/fig06_evidence_map.pdf}
\caption{Evidence flow from generation studies to read-only validation, constructed-state analysis, and historical-pool selection. The four empirical streams answer complementary questions and are not pooled into one framework accuracy.}
\label{fig:evidence-map}
\end{figure}

\subsection{Limitations and Next Evaluation}

GridDB contains eight tables and 98 synthetic rows, and its evaluation questions were visible during development. The question groups used for resampling are dependence proxies, not sampled deployment clusters. BIRD adds public database diversity but no independent validation of power-grid business semantics. The findings therefore describe the retained cases and model snapshots rather than a deployment population.

Chronology imposes a second limit. Although selection is sealed before offline scoring, same-item derived outcomes had been accessed during earlier development. The reconciled evaluator, global-order enumeration, and role diagnostics are post-result analyses. Their transparency does not convert them into confirmation on unseen data. Likewise, later executor hardening cannot be credited to the earlier ledgers.

The formal properties are intentionally modest. They characterize a finite selector under fixed inputs and correct implementation; they do not establish learned calibration, general SQL equivalence, or system security beyond the stated threat model. No expert user study measures review time, usability, or unsafe omission. Existing machine-adjudicated RTS-GMLC, SimBench, and NERC-derived question--SQL resources cannot supply that missing human evidence.

The next evaluation should freeze schema and value retrieval, candidate budgets, semantic evidence, tie handling, and abstention rules before opening unseen grid schemas and independently reviewed questions. It should compare direct generation, staged single-candidate processing, and multi-candidate selection under matched physical calls and tokens. Reported outcomes should include reference agreement, expert-semantic errors, correct-candidate loss at gates, coverage, and total evaluation cost. These are proposed requirements, not experiments completed in this revision.

\section{Conclusions}

\masql{} separates candidate production, eligibility checks, and final selection in an auditable power-grid Text-to-SQL workflow. Its five role interfaces and deterministic controller expose the evidence behind each decision. The formal analysis establishes limited properties of that rule: selection stays within the eligible pool, its accuracy is bounded by available correct candidates, and unresolved top ties delegate the decision to candidate order. None of these properties implies semantic correctness.

The 180-question historical pool separates the remaining errors into actionable categories: 42 questions contain no correct candidate, no evaluator-correct slot is lost at eligibility, six questions lose the correct candidates at scoring, and 33 or 32 lose a correct top candidate through the original tie order. Validation-aware and complete-witness selection therefore match 99 and 100 references, below the strongest fixed source at 129. Mixed-correctness top sets occur on 51 questions, not on every one of the 130 slot-tied questions. These are finite-case findings about the recorded rule; the supporting BIRD generation results do not establish their prevalence across databases.

The design lesson is to diagnose the decision stage before adding complexity. A missing correct candidate requires a change upstream of selection; a correct candidate below the top score requires better discrimination; a mixed-correctness top set requires additional evidence or a justified ambiguity policy. More eligibility checks alone need not repair either of the latter failures. Whether this decomposition identifies the same bottlenecks in other settings must be tested on independently reviewed queries and new candidate pools under a frozen, budget-matched protocol.

\supplementary{The Supplementary Materials include software and executor tests, protocol and version records, complete numerical tables, figure provenance, rights information, and reproducibility instructions.}
\authorcontributions{Conceptualization, B.L. and Y.Y.; methodology, B.L.; software, B.L.; validation, B.L. and C.S.; formal analysis, B.L. and C.S.; investigation, B.L. and C.S.; resources, Y.Y.; data curation, B.L. and C.S.; writing---original draft preparation, B.L.; writing---review and editing, C.S. and Y.Y.; visualization, B.L.; supervision, Y.Y.; project administration, Y.Y.; funding acquisition, Y.Y. All authors have read and agreed to the published version of the manuscript.}
\funding{This research was funded by the Science and Technology Research Project of State Grid Fujian Electric Power Co., Ltd., grant number 521300250006. The funder had no role in the study design; in the collection, analysis, or interpretation of data; in the writing of the manuscript; or in the decision to publish the results.}
\institutionalreview{Not applicable. This study did not involve human participants or animals.}
\informedconsent{Not applicable.}
\dataavailability{The manuscript-bound, rights-safe technical package can be inspected in the public \nolinkurl{powergrid_benchmark} repository under the frozen tag \nolinkurl{cmc-2026-08-24-v3} via the \href{https://github.com/gaoxingkele/powergrid_benchmark/tree/cmc-2026-08-24-v3/paper_projects/CMC/MA-SQLGrid}{frozen MA-SQLGrid package}. It includes the numerical supplement, code, tests, derived evidence, and reproducibility records underlying the historical empirical results; the present Information rewrite is not contained in that historical tag. BIRD Mini-Dev remains available from its public provider. Owing to third-party licensing restrictions, raw GridDB and BIRD databases and restricted source records are not redistributed, and source-dependent RTS-GMLC or SimBench artifacts may be provided for editorial or peer-review verification only when permitted by the applicable licenses. The historical \texttt{ma-sqlgrid} repository is not the manuscript-bound release for this version.}
\acknowledgments{During preparation of this manuscript, the authors used the OpenAI Codex agent interface for the activities described in Materials and Methods, including manuscript preparation and local code and build review. The exact backend model identifier and version were not retained for every session. The authors reviewed and edited the assisted work and take full responsibility for the analyses, interpretations, and content of this publication.}
\conflictsofinterest{The authors declare no conflicts of interest.}

\appendix
\section{Historical Question-Level Statistical Reconciliation}
\label{sec:historical-question-statistics}

Table~\ref{tab:historical-question-statistics} preserves the original post-result numerical reconciliation. Its resampling unit is the individual question, unlike the grouped analyses of the generation studies. Dependence among questions sharing SQL structure is not addressed by these intervals or discordant-sign tests. They are retained for provenance, not used to support generalization or the diagnostic contribution. The main text relies on exhaustive finite-pool counts and the stage partition.

\begin{table}[htbp]
\caption{Archived question-level reconciliation, without a cross-template or population interpretation. Intervals are paired question-composition summaries; Holm adjustment covers nine unique C000 comparisons.}
\label{tab:historical-question-statistics}\centering\scriptsize
\resizebox{\textwidth}{!}{%
\begin{tabular}{lrrrrrr}
\toprule
Method & Correct/180 & Accuracy & Difference & 95\% comp. interval & Rescue/harm & Holm $p$\\
\midrule
C000 / Qwen F00 & 76/180 & 0.4222 & 0 & [0, 0] & 0/0 & --\\
Qwen F01 & 129/180 & 0.7167 & +0.2944 & [+0.2222, +0.3667] & 56/3 & $<0.000001$\\
Qwen F10 & 78/180 & 0.4333 & +0.0111 & [-0.0111, +0.0333] & 3/1 & 1.000000\\
Qwen F11 & 108/180 & 0.6000 & +0.1778 & [+0.1000, +0.2556] & 46/14 & 0.000211\\
Granite F00 & 77/180 & 0.4278 & +0.0056 & [-0.0278, +0.0444] & 6/5 & 1.000000\\
Granite F01 & 100/180 & 0.5556 & +0.1333 & [+0.0500, +0.2222] & 45/21 & 0.017090\\
Granite F10 & 74/180 & 0.4111 & -0.0111 & [-0.0500, +0.0278] & 6/8 & 1.000000\\
Granite F11 & 108/180 & 0.6000 & +0.1778 & [+0.1056, +0.2500] & 42/10 & 0.000054\\
Validation-only selector & 99/180 & 0.5500 & +0.1278 & [+0.0778, +0.1833] & 25/2 & 0.000040\\
Complete-witness selector & 100/180 & 0.5556 & +0.1333 & [+0.0778, +0.1889] & 26/2 & 0.000024\\
\bottomrule
\end{tabular}}
\end{table}

\begin{adjustwidth}{-\extralength}{0cm}
\reftitle{References}
\bibliography{references_verified,references_information}
\PublishersNote{}
\end{adjustwidth}
\end{document}
